\SourceLeaf{112}
\expandafter\def\csname HistoricalNoteText12\endcsname{Longitudo acuum in his experimentis adhibitarum est circiter $0^{\mathrm m}3$; si terminum $R^{-7}$ in calculis respicere periclitati essemus, certitudo minuta potius quam aucta fuisset.}
\Needspace{13\baselineskip}
\begin{center}
\textit{Modus secundus.}\par
Acus secundae centrum in meridiano magnetico, axis huic normalis.
\begin{tabular}{c|cc|cc}
Deflexio&\multicolumn{4}{c}{Situs acus}\\
&&&centrum versus&polus borealis versus\\\hline
$u=u'$&$\psi=0$&$U=270^0$&boream&occidentem\\
$u=u''$&$\psi=0$&$U=270^0$&boream&orientem\\
$u=u'''$&$\psi=180^0$&$U=90^0$&austrum&occidentem\\
$u=u''''$&$\psi=180^0$&$U=90^0$&austrum&orientem
\end{tabular}
\end{center}

Statuendo dein $\frac14(u'-u''+u'''-u'''')=v$, atque
\[
\operatorname{tang}v=LR^{-(n+1)}+L'R^{-(n+3)}+L''R^{-(n+5)}+\text{etc.}
\]
erit
\[
\begin{aligned}
\text{pro modo priori}\qquad L&=\frac{nmM}{mT+\theta}\\
\text{pro modo posteriori}\qquad L&=\frac{mM}{mT+\theta}
\end{aligned}
\]
\Article{20}
E theoria eliminationis facile colligitur, calculum, propter inevitabiles observationum errores, eo magis incertum fieri. quo plures coëfficientes per eliminationem determinare oporteat. Hanc ob caussam modus in art. $18$, II. praescriptus magni aestimandus est, quod coëfficientes potestatum $R^{-(n+2)},R^{-(n+4)}$ supprimit. In casu perfectae symmetriae quidem hi termini iam per se exciderent, sed parum tutum esset, illi fidem habere. Ceterum parvula a symmetria aberratio longe minoris momenti esset in modo primo quam in secundo, et si in illo saltem cavetur, ut punctum $h'$. a quo distantiae numerantur, sit satis exacte in meridiano magnetico per $h$ transeunte, vix differentia sensibilis inter $u'-u''$ atque $u'''-u''''$ se manifestabit. Sed hoc secus se habet in modo secundo, praesertim si apparatus suspensionem excentricam postulat. Per hunc modum, quoties spatium non permittit observationes ab utraque parte, semper praecisionem multo minorem assequeris. Praeterea modus primus eo quoque nomine praeferendus est, quod, quum in casu naturae sit $n=2$, duplo maiorem valorem ipsius $L$ producit, quam secundus. Ceterum si in modo secundo terminum a $R^{-(n+2)}$ pendentem, in casu suspensionis excentricae, quantum licet exterminare studemus,
\SourceLeaf{113}
punctum $h'$ ita eligendum est, ut centrum acus (pro $u=u^0$) sit medium inter $h$ et $h'$: calculum tamen, qui hoc docuit, brevitatis caussa hic supprimere oportet.

\Article{21}
In calculis praecedentibus exponentem $n$ indeterminatum liquimus: diebus Iunii $24$—$28$, $1832$ duas series experimentorum exsequuti sumus, ad tantas distantias, quantas spatium permisit, extensas, per quas, quemnam valorem natura postulet, evidentissime apparebit. In prima serie acus secunda (ad modum primum art. $19$) in recta ad meridianum magneticum normali, in secunda centrum acus in ipso meridiano collocabatur. Ecce conspectum summae horum experimentorum, ubi distantiae $R$ in partibus metri expressae, valoresque anguli $\frac14(u-u'+u''-u''')$ pro prima serie per $v$, pro secunda per $v'$ denotati sunt.
\begin{center}
\begin{tabular}{c|rrr|rrr}
$R$&\multicolumn{3}{c|}{$v$}&\multicolumn{3}{c}{$v'$}\\\hline
$1^{\mathrm m}1$&&&&$1^0$&$57'$&$24''8$\\
$1{,}2$&&&&$1$&$29$&$40{,}5$\\
$1{,}3$&$2^0$&$13'$&$51''2$&$1$&$10$&$19{,}3$\\
$1{,}4$&$1$&$47$&$28{,}6$&$0$&$55$&$58{,}9$\\
$1{,}5$&$1$&$27$&$19{,}1$&$0$&$45$&$14{,}3$\\
$1{,}6$&$1$&$12$&$7{,}6$&$0$&$37$&$12{,}2$\\
$1{,}7$&$1$&$0$&$9{,}9$&$0$&$30$&$57{,}9$\\
$1{,}8$&$0$&$50$&$52{,}5$&$0$&$25$&$59{,}5$\\
$1{,}9$&$0$&$43$&$21{,}8$&$0$&$22$&$9{,}2$\\
$2{,}0$&$0$&$37$&$16{,}2$&$0$&$19$&$1{,}6$\\
$2{,}1$&$0$&$32$&$4{,}6$&$0$&$16$&$24{,}7$\\
$2{,}5$&$0$&$18$&$51{,}9$&$0$&$9$&$36{,}1$\\
$3{,}0$&$0$&$11$&$0{,}7$&$0$&$5$&$33{,}7$\\
$3{,}5$&$0$&$6$&$56{,}9$&$0$&$3$&$28{,}9$\\
$4{,}0$&$0$&$4$&$35{,}9$&$0$&$2$&$22{,}2$
\end{tabular}
\end{center}

Hi numeri vel obiter inspecti monstrant, pro valoribus maioribus tum numeros secundae columnae proxime duplo maiores esse numeris tertiae, tum numeros utriusque columnae proxime in ratione inversa cubi distantiarum; ita ut de veritate valoris $n=2$ nullum dubium remanere possit. Quo magis adhuc haec
\SourceLeaf{114}
lex in singulis experimentis confirmaretur, omnes numeros per methodum quadratorum minimorum tractavimus, unde valores coëfficientium sequentes prodierunt:
\[
\begin{aligned}
\operatorname{tang}v&=0{,}086870R^{-3}-0{,}002185R^{-5}\\
\operatorname{tang}v'&=0{,}043435R^{-3}+0{,}002449R^{-5}
\end{aligned}
\]

Ecce conspectum comparationis valorum per has formulas computatorum cum observatis:
\begin{center}
\textit{Valores computati.}\par
\setlength{\tabcolsep}{4pt}
\begin{tabular}{c||rrr|r||rrr|r}
$R$&\multicolumn{3}{c|}{$v$}&differentia&\multicolumn{3}{c|}{$v'$}&differentia\\\hline
$1^{\mathrm m}1$&&&&&$1^0$&$57'$&$22''0$&$+2''8$\\
$1{,}2$&&&&&$1$&$29$&$46{,}5$&$-6{,}0$\\
$1{,}3$&$2^0$&$13'$&$50''4$&$+0''8$&$1$&$10$&$13{,}3$&$+6{,}0$\\
$1{,}4$&$1$&$47$&$24{,}1$&$+4{,}5$&$0$&$55$&$58{,}7$&$+0{,}2$\\
$1{,}5$&$1$&$27$&$28{,}7$&$-9{,}6$&$0$&$45$&$20{,}9$&$-6{,}6$\\
$1{,}6$&$1$&$12$&$10{,}9$&$-3{,}3$&$0$&$37$&$15{,}4$&$-3{,}2$\\
$1{,}7$&$1$&$0$&$14{,}9$&$-5.0$&$0$&$30$&$59{,}1$&$-1{,}2$\\
$1{,}8$&$0$&$50$&$48{,}3$&$+4{,}2$&$0$&$26$&$2{,}9$&$-3{,}4$\\
$1{,}9$&$0$&$43$&$14{,}0$&$+7{,}8$&$0$&$22$&$6{,}6$&$+2{,}6$\\
$2{,}0$&$0$&$37$&$5{,}6$&$+10{,}6$&$0$&$18$&$55{,}7$&$+5{,}9$\\
$2{,}1$&$0$&$32$&$3{,}7$&$+0{,}9$&$0$&$16$&$19{,}8$&$+4{,}9$\\
$2{,}5$&$0$&$19$&$2{,}1$&$-10{,}2$&$0$&$9$&$38{,}6$&$-2{,}5$\\
$3{,}0$&$0$&$11$&$1{,}8$&$-1{,}1$&$0$&$5$&$33{,}9$&$-0{,}2$\\
$3{,}5$&$0$&$6$&$57{,}1$&$-0{,}2$&$0$&$3$&$29{,}8$&$-1{,}0$\\
$4{,}0$&$0$&$4$&$39.6$&$-3{,}7$&$0$&$2$&$20{,}5$&$+1{,}7$
\end{tabular}
\end{center}

\Article{22}
Experimenta praecedentia eum potissimum in finem suscepta fuerunt, ut lex actionis magneticae ultra omne dubium eveheretur, porro, ut quot terminos seriei respicere conveniat, quantamque praecisionem ferant experimenta, appareret. Docuerunt, si ad distantias minores quadruplo longitudinis acuum non de\SourceLeaf{115}scendamus, duos seriei terminos sufficere\HistoricalNote{12}. Ceterum differentiae, quas calculus prodidit, neutiquam pure pro erroribus observationum haberi debent: plures enim cautelae, a quarum usu harmoniam adhuc maiorem sperare licet, tunc temporis nondum praeparatae erant. Huc referendae sunt correctiones propter variabilitatem horariam intensitatis magnetismi terrestris, cuius rationem habere oportet adiumento alius acus comparativae ad instar methodi, de qua in art. $10$, II. loquuti sumus. Quo tamen valor magnetismi terrestris, quatenus ex \textit{his} experimentis derivari potest, cognoscatur, summam reliquorum experimentorum huc spectantium adiicimus.

Valor fractionis $\frac{\theta}{Tm}$ pro acu prima et filo, a quo pendebat, erutus est per methodum in art. $8$ traditam $=\frac{1}{251{,}96}$. Hinc itaque fit
\[
\frac MT=0{,}0436074
\]
Huic numero subest metrum tamquam unitas distantiarum. Si pro unitate millimetrum adoptare malumus, iste numerus per cubum millenarii multiplicandus est, ita ut habeatur
\[
\frac MT=43607400
\]

Pro acu secunda per experimenta d. $28$. Junii instituta iisque prorsus similia, quae pro alia acu in art. $11$ tractavimus, prodiit, dum millimetrum, milligramma et minutum secundum temporis solaris medii pro unitatibus accipiebantur,
\[
TM=135457900
\]
atque hinc, eliminata quantitate $M$
\[
T=1{,}7625
\]
\Article{23}
Quoties experimenta eum in finem instituuntur, ut valor absolutus magnetismi terrestris $T$ determinetur, magni momenti est, curare, ut ipsorum complexus intra modicum tempus absolvatur, ne mutatio sensibilis status magnetici acuum ad illa adhibitarum metuenda sit. Conveniet itaque in observandis deflexionibus acus mobilis solum modum primum art. $20$ sequi, adhibitis tantummodo
\SourceLeaf{116}
duabus distantiis diversis apte electis, siquidem duo seriei termini sufficiunt. E pluribus applicationibus huius methodi hic unam tamquam exemplum eligimus, et quidem eam, cui cura maxime scrupulosa impensa est, distantiis microscopica praecisione mensuratis.

Experimenta instituta sunt $1832$ Sept. $18$, in duobus apparatibus, quos per literas $A,B$, tribus acubus, quas per numeros $1,2,3$ distinguemus. Acus $1,2$ sunt eaedem, quae in art. $11$ prima et secunda vocabantur. Experimenta ad duo capita discedunt.

Primo observatae sunt oscillationes simultaneae acus $1$ in apparatu $A$, acusque $2$ in apparatu $B$. Tempus unius oscillationis, ad amplitudinem infinite parvam reductum prodiit
\begin{center}
\begin{tabular}{lcl}
pro acu $1$&.\ .\ .\ .\ .\ .&$15''22450$\\
pro acu $2$&.\ .\ .\ .\ .\ .&$17{,}29995$
\end{tabular}
\end{center}
illud ex $305$, hoc ex $264$ oscillationibus conclusum.

Dein acus $3$ in apparatu $A$ suspensa, acus $1$ autem in recta ad meridianum magneticum normali tum versus orientem tum versus occidentem, et utrimque duplici modo collocata, deflexioque acus $3$ pro singulis positionibus acus $1$ observata est. Haec experimenta, pro duabus distantiis diversis $R$ repetita prodiderunt valores sequentes anguli $v$ perinde intelligendi ut in artt. $19$, $21$
\[
\begin{aligned}
R&=1^{\mathrm m}2,&\qquad v&=3^0\ 42'\ 19''4\\
R'&=1{,}6.&v'&=1\ 34\ 19{,}3
\end{aligned}
\]
Inter haec quoque experimenta oscillationes acus $2$ in apparatu $B$ observatae sunt: tempori medio respondet valor unius oscillationis infinite parvae ex $414$ oscillationibus conclusus $=17''29484$.

Tempora observabantur ad chronometrum, cuius retardatio diurna $=14''24$.

Denotantibus $M,m$ momenta magnetismi liberi pro acu $1$ et $3$, $\theta$ constantem torsionis fili in apparatu $A$, dum acum $1$ vel $3$ (quarum pondus fere idem est) ferebat, habemus
\[
\frac{\theta}{TM}=\frac1{597{,}4}
\]
uti in art. $11$
\[
\frac{\theta}{Tm}=\frac1{721{,}6}
\]
quippe acus $3$ fortiori magnetismo imbuta erat, quam acus $1$.
